\documentclass[10pt,a4paper]{article}
\usepackage{amsmath,amssymb,amsthm,graphicx}
\usepackage{hyperref}
\usepackage{cleveref}
\usepackage{amssymb,amsmath,amsthm}
\theoremstyle{plain}
\newtheorem{theorem}{Theorem}
\newtheorem{lemma}[theorem]{Lemma}
\newtheorem{observation}[theorem]{Observation}
\theoremstyle{definition}
\newtheorem{remark}[theorem]{Remark}
\title{A solution of bounded dynamic range yields a linearly ordered colouring with logarithmically many colours}
\author{Johan Hastad, Bjorn Martinsson, Tamio-Vesa Nakajima, Stanislav Zivny}
\date{arxiv:2404.19556v6; Creative Commons Attribution 4.0 International}
\begin{document}
\maketitle
\noindent\textit{Editorial note.} The counted statement is the author's labelled environment at source0.tex lines 266-271 and its proof environment opens at line 273; the pairing is by adjacency and the retained proof range is the whole environment.\par\medskip
\noindent\textit{Editorial note.} This source is newly approved for this line and its preamble is built from the author's own package line and theorem-environment family (source lines 5-23); the font, language and layout packages of the source are not retained because the retained material does not use them.\par\medskip
\noindent\textit{Editorial note.} Metadata provenance: the pinned source carries a licence evidence file recording source\_version 2404.19556v6, CC BY 4.0, and a separate source record file; the card's source\_version is the value recorded in the licence evidence.\par\medskip
\noindent\textit{Editorial note.} The source carries no label-number dump and no counter values are imposed, so the printed numbering is the wrapper's default; identity is the author label.\par\medskip
\noindent\textit{Editorial note.} r002 note: r001 was refused by the assembler because the card carried an empty author list; r002 fills in the authors and the source title recorded in the source record of the pinned copy and changes nothing else.\par\medskip
\noindent\textit{Editorial note.} r003 note: the pre-publication validator requires either status=supported in the licence evidence or a retained-licence form with no status at all. This source's evidence file carries no acceptance record and none may be written on anyone's behalf, so r003 keeps no status and records a retained source notice (source\_notice), quoting the author's own CC BY declaration on lines 47-48 of the manuscript with its path form in a task-directory evidence copy that the assembler carries into the package: the manuscript declares CC BY at lines 47-48. The original keys of the source evidence are retained unchanged, no new licence claim is made, and no acceptance record is written.\par\medskip
\noindent\textit{Editorial scope.} One original result of the pinned author manuscript, counted once: a linearly ordered colouring with logarithmically many colours from a solution of bounded dynamic range (author label lemma:unbal).\par\medskip
\noindent\textit{Allowed standard prerequisites (editorial).}\begin{itemize}
\item Homogeneous linear system over the reals. The coordinates of a solution are real numbers, with M the largest absolute value and m the smallest positive absolute value.. Their ratio controls the number of colours.
\item Linearly ordered colouring. A colouring in which the smallest colour appearing once in an edge is unique.. It is what the construction produces.
\item Logarithmic colour classes. The colours are assigned by comparing the absolute value of a coordinate against powers of two, with even colours used for positive coordinates and odd ones for negative coordinates.. It is the construction whose correctness has to be checked.
\end{itemize}\medskip
\begin{lemma}\label{lemma:unbal}
Suppose we have a solution, $u$, to the homogeneous system where $M$
is the maximal value of the absolute value of a coordinate
and $m > 0$ is the minimal absolute value.  Then, we can
LO colour with $2+\log_2{(M/m)}$ colours.
\end{lemma}

\begin{proof} For notational convenience let us instead require that the
minimal colour in each edge should be unique. We can simply reverse
the order of the colours at the end.  By scaling we can assume $M=1$.
We use even colours for positive coordinates and we
give $x$ the colour $2\ell$ if $v_x$ is at most $2^{-(2\ell-1)}$
and strictly larger than $2^{-(2\ell+1)}$.
For negative coordinates we use $2\ell+1$ as
the colour if $v_x$ is between $-2^{-2\ell}$ (inclusive)
and $-2^{-(2\ell+2)}$ (non-inclusive).
Let us verify that this gives a correct
colouring.

Take an edge $(x,y,z)$ and suppose both $x$ and $y$
get the same colour $2\ell$.   Then by the linear equation of the edge $v_z < -2^{-2\ell}$ and thus $z$ has a colour below $2\ell$.  The case
of two vertices of odd colour is similar and as the bound on the number of colours is immediate, the lemma follows.
\end{proof}

\end{document}
